% GENERATED FILE. Edit canonical lessons, then run npm run generate:books.
% canonical-source-hash: fnv1a64:e5d88f839538e887
% canonical-lessons: GE-C420-stornieren, GE-C420-stornierung, GE-C420-umbuchen, GE-C420-erstatten, GE-C420-erstattung

\chapter{Flug gestrichen}
\label{ch:ge-flug-gestrichen}
\hlchaptermodality{420}

\begin{chapteropening}
\textbf{By the end of this chapter:} \emph{I can deal with a cancelled booking: cancelling, the cancellation, rebooking, refunding and the refund.}

Complete the last lesson of chapter 420: I can deal with a cancelled booking: cancelling, the cancellation, rebooking, refunding and the refund.
\end{chapteropening}

\section[stornieren]{stornieren — to cancel (a booking)}

\label{lesson:GE-C420-stornieren}



\begin{quote}
Before the new one: say the German for to be stuck, then the German for to make an emergency landing.
\end{quote}

\subsection*{You'll want to know: stornieren}
\textbf{stornieren} — \textquotedblleft{}to cancel (a booking)\textquotedblright{}.

\begin{quote}
Er storniert das Hotel, weil er krank ist.
\end{quote}

\begin{quote}
\emph{He cancels the hotel because he is ill.}
\end{quote}

Past participle without \emph{ge-}: \emph{ich habe den Flug storniert}.

\begin{cousinweb}
\emph{stornieren} was borrowed from Italian \emph{stornare} 'to turn back, to reverse an entry in the books' (EtymDB), a word from bookkeeping.
\end{cousinweb}

\begin{grammarlens}[title={What you've built}]
One more word for telling, explaining and describing at B1.
\end{grammarlens}

\subsection*{Your turn}
\begin{itemize}
\raggedright
  \item \emph{Say it:} \emph{stornieren}
  \item \emph{Say it:} \emph{stornieren}, once more
  \item \emph{Say it:} a sentence of your own with \emph{stornieren}
  \item \emph{Recall:} say the German for to be stuck, then the German for to make an emergency landing, then say \emph{stornieren} again
\end{itemize}

\begin{tcolorbox}[breakable,colback=teal!4,colframe=teal!35!black,title={Before you move on}]
What does stornieren mean? (\textquotedblleft{}To cancel (a booking)\textquotedblright{}.) Say it once more.
\end{tcolorbox}
\section[die Stornierung]{die Stornierung — the cancellation}

\label{lesson:GE-C420-stornierung}



\begin{quote}
Before the new one: say the German for to make an emergency landing, then the German for to cancel.
\end{quote}

\subsection*{You'll want to know: die Stornierung}
\textbf{die Stornierung} — \textquotedblleft{}the cancellation\textquotedblright{}.

\begin{quote}
Die Stornierung kostet zwanzig Euro.
\end{quote}

\begin{quote}
\emph{The cancellation costs twenty euros.}
\end{quote}

Plural: \emph{die Stornierungen}. Also \emph{das Storno}.

\begin{cousinweb}
\emph{Stornierung} is built on \emph{stornieren} with the ending \emph{-ung}; the verb was borrowed from Italian \emph{stornare} (EtymDB).
\end{cousinweb}

\begin{grammarlens}[title={What you've built}]
One more word for telling, explaining and describing at B1.
\end{grammarlens}

\subsection*{Your turn}
\begin{itemize}
\raggedright
  \item \emph{Say it:} \emph{die Stornierung}
  \item \emph{Say it:} \emph{die Stornierung}, once more
  \item \emph{Say it:} a sentence of your own with \emph{die Stornierung}
  \item \emph{Recall:} say the German for to make an emergency landing, then the German for to cancel, then say \emph{die Stornierung} again
\end{itemize}

\begin{tcolorbox}[breakable,colback=teal!4,colframe=teal!35!black,title={Before you move on}]
What does die Stornierung mean? (\textquotedblleft{}The cancellation\textquotedblright{}.) Say it once more.
\end{tcolorbox}
\section[umbuchen]{umbuchen — to rebook, to change a booking}

\label{lesson:GE-C420-umbuchen}



\begin{quote}
Before the new one: say the German for to cancel, then the German for the cancellation.
\end{quote}

\subsection*{You'll want to know: umbuchen}
\textbf{umbuchen} — \textquotedblleft{}to rebook, to change a booking\textquotedblright{}.

\begin{quote}
Die Fluggesellschaft hat uns auf einen späteren Flug umgebucht.
\end{quote}

\begin{quote}
\emph{The airline has rebooked us onto a later flight.}
\end{quote}

Separable: \emph{ich buche um}, \emph{ich habe umgebucht}.

\begin{cousinweb}
\emph{umbuchen} is built from \emph{um} 'over, round' and \emph{buchen} 'to book', first 'to enter in a book', from \emph{Buch} 'book' (EtymDB).
\end{cousinweb}

\begin{grammarlens}[title={What you've built}]
One more word for telling, explaining and describing at B1.
\end{grammarlens}

\subsection*{Your turn}
\begin{itemize}
\raggedright
  \item \emph{Say it:} \emph{umbuchen}
  \item \emph{Say it:} \emph{umbuchen}, once more
  \item \emph{Say it:} a sentence of your own with \emph{umbuchen}
  \item \emph{Recall:} say the German for to cancel, then the German for the cancellation, then say \emph{umbuchen} again
\end{itemize}

\begin{tcolorbox}[breakable,colback=teal!4,colframe=teal!35!black,title={Before you move on}]
What does umbuchen mean? (\textquotedblleft{}To rebook, to change a booking\textquotedblright{}.) Say it once more.
\end{tcolorbox}
\section[erstatten]{erstatten — to refund, to pay back}

\label{lesson:GE-C420-erstatten}



\begin{quote}
Before the new one: say the German for the cancellation, then the German for to rebook.
\end{quote}

\subsection*{You'll want to know: erstatten}
\textbf{erstatten} — \textquotedblleft{}to refund, to pay back\textquotedblright{}.

\begin{quote}
Die Firma erstattet uns die Kosten für das Hotel.
\end{quote}

\begin{quote}
\emph{The company refunds us the cost of the hotel.}
\end{quote}

Formal. \emph{Anzeige erstatten} — to report a crime to the police.

\begin{cousinweb}
\emph{erstatten} is built from \emph{er-} and the old verb \emph{statten} 'to put in place', from \emph{Statt} 'place', an inherited Germanic word (EtymDB lists \emph{erstatten}): to put money back where it was.
\end{cousinweb}

\begin{grammarlens}[title={What you've built}]
One more word for telling, explaining and describing at B1.
\end{grammarlens}

\subsection*{Your turn}
\begin{itemize}
\raggedright
  \item \emph{Say it:} \emph{erstatten}
  \item \emph{Say it:} \emph{erstatten}, once more
  \item \emph{Say it:} a sentence of your own with \emph{erstatten}
  \item \emph{Recall:} say the German for the cancellation, then the German for to rebook, then say \emph{erstatten} again
\end{itemize}

\begin{tcolorbox}[breakable,colback=teal!4,colframe=teal!35!black,title={Before you move on}]
What does erstatten mean? (\textquotedblleft{}To refund, to pay back\textquotedblright{}.) Say it once more.
\end{tcolorbox}
\section[die Erstattung]{die Erstattung — the refund}

\label{lesson:GE-C420-erstattung}



\begin{quote}
Before the new one: say the German for to rebook, then the German for to refund.
\end{quote}

\subsection*{You'll want to know: die Erstattung}
\textbf{die Erstattung} — \textquotedblleft{}the refund\textquotedblright{}.

\begin{quote}
Die Erstattung kommt nach zwei Wochen auf das Konto.
\end{quote}

\begin{quote}
\emph{The refund reaches the account after two weeks.}
\end{quote}

Plural: \emph{die Erstattungen}. Also \emph{die Rückerstattung}.

\begin{cousinweb}
\emph{Erstattung} is built on \emph{erstatten} 'to refund' with \emph{-ung}; the verb comes from \emph{er-} and \emph{statten} 'to put in place', from \emph{Statt} 'place'.
\end{cousinweb}

\begin{grammarlens}[title={What you've built}]
One more word for telling, explaining and describing at B1.
\end{grammarlens}

\subsection*{Your turn}
\begin{itemize}
\raggedright
  \item \emph{Say it:} \emph{die Erstattung}
  \item \emph{Say it:} \emph{die Erstattung}, once more
  \item \emph{Say it:} a sentence of your own with \emph{die Erstattung}
  \item \emph{Recall:} say the German for to rebook, then the German for to refund, then say \emph{die Erstattung} again
\end{itemize}

\begin{tcolorbox}[breakable,colback=teal!4,colframe=teal!35!black,title={Before you move on}]
What does die Erstattung mean? (\textquotedblleft{}The refund\textquotedblright{}.) Say it once more.
\end{tcolorbox}
